\documentclass[12pt]{article}
\usepackage[T1]{fontenc}
\usepackage[utf8]{inputenc}
\usepackage{lmodern}
\usepackage{textcomp}
\usepackage{enumitem}
\usepackage{geometry}
\usepackage{graphicx}
\usepackage{amsmath, amssymb}
\geometry{margin=1in}
\begin{document}
\begin{center}
Physics - Undergraduate Level Assessment 11 \
Total Marks: 40 \
Answer all questions.
\end{center}

\section*{Multiple Choice (10 marks)}
\begin{enumerate}[label=\arabic*.]
\item The sign of the charge carriers in a doped semiconductor can be deduced by measuring which of the following properties?
\begin{enumerate}[label=(\alph*)]
    \item Magnetic susceptibility
    \item Hall coefficient
    \item Electrical resistivity
    \item Thermal conductivity
\end{enumerate}
\item A single-electron atom has the electron in the l = 2 state. The number of allowed values of the quantum number m\_l is
\begin{enumerate}[label=(\alph*)]
    \item 5
    \item 4
    \item 3
    \item 2
\end{enumerate}
\item A rod measures 1.00 m in its rest system. How fast must an observer move parallel to the rod to measure its length to be 0.80 m?
\begin{enumerate}[label=(\alph*)]
    \item 0.50 c
    \item 0.60 c
    \item 0.70 c
    \item 0.80 c
\end{enumerate}
\item De Broglie hypothesized that the linear momentum and wavelength of a free massive particle are related by which of the following constants?
\begin{enumerate}[label=(\alph*)]
    \item Planck's constant
    \item Boltzmann's constant
    \item The Rydberg constant
    \item The speed of light
\end{enumerate}
\item A refracting telescope consists of two converging lenses separated by 100 cm. The eye- piece lens has a focal length of 20 cm. The angular magnification of the telescope is
\begin{enumerate}[label=(\alph*)]
    \item 4
    \item 5
    \item 6
    \item 20
\end{enumerate}
\end{enumerate}

\section*{True / False (10 marks)}
\textit{Answer True or False}
\begin{enumerate}[label=\arabic*.]
\setcounter{enumi}{5}
\item True or False: Fermions have antisymmetric wave functions and obey the Pauli exclusion principle.
\item True or False: A particle moving at v=0.60 c will travel 288 m before decaying in the lab frame.
\item True or False: The angular magnification of the refracting telescope is 4, given the eye-piece lens has a focal length of 20 cm.
\item True or False: The mean kinetic energy of conduction electrons in metals is high due to their many degrees of freedom compared to atoms.
\item True or False: If the absolute temperature of a blackbody is increased by a factor of 3, the energy radiated per second per unit area increases by a factor of 81.
\end{enumerate}

\section*{Long Form Response (20 marks)}
\textit{Answer each question with a clear and complete explanation.}
\begin{enumerate}[label=\arabic*.]
\setcounter{enumi}{10}
\item Discuss the total spin quantum number of the electrons in the ground state of neutral nitrogen (Z = 7), explaining how you arrive at your conclusion and its significance in atomic structure.
\item Discuss the electron configuration of an atom with filled n = 1 and n = 2 energy levels, and analyze how this configuration determines the total number of electrons in the atom.
\end{enumerate}

\vfill
\noindent\textit{End of Paper}
\end{document}